%% parts/p04/decision-problems-and-arithmetical-hierarchy.tex
%% Migrates v1.1 lines 8248--8897.  The repaired R-002/R-003/R-006
%% definitions and dependency order are preserved; proofs are in the parallel
%% annex fragment.

\chapter[Decision Problems and Hierarchy]{Decision Problems and the Arithmetical Hierarchy}
\label{sec:p04-decision-hierarchy}

This chapter fixes one effective coding of first-order formulas and finite
derivations and uses the machine numbering of the preceding chapter.  The
coding is part of each computability classification: changing a code is
harmless only after a computable translation has been supplied.  Standard
references are \cite{Church1936,Turing1937,Shoenfield1967,
BoolosBurgessJeffrey2007}.

\section{Logical decision problems}
\label{sec:p04-logical-decision-problems}

\begin{hbtheorem}{thm:p04-fol-validity-undecidable}{cited=Church1936}{Negative solution to the classical Entscheidungsproblem}
Under the fixed effective coding, no Turing machine decides the set of codes
of logically valid first-order sentences.  Turing's parallel formulation is
recorded in \cite{Turing1937}; the claim concerns the general decision problem,
not every fixed signature or fragment.
\end{hbtheorem}

\begin{hbcorollary}{thm:p04-fol-validity-ce}{proved=annex:p04-hierarchy.1}{Enumerability without decidability}
Under the fixed effective coding, the set of logically valid first-order
sentences is computably enumerable but not decidable.
\end{hbcorollary}

\begin{hbwarning}{def:p04-restricted-signatures}{Restricted signatures}
The negative solution concerns unrestricted first-order validity.  The monadic
fragment without function symbols is decidable, so undecidability cannot be
transferred to an individual signature or fragment without a theorem that
controls that restriction.
\end{hbwarning}

\begin{hbtheorem}{thm:p04-trakhtenbrot}{cited=Trakhtenbrot1950}{Trakhtenbrot's theorem}
For every finite relational signature containing a binary relation symbol, the
set of sentences having a finite model is computably enumerable but
undecidable, while the set of sentences true in every finite structure is not
computably enumerable.
\end{hbtheorem}

Unrestricted validities are computably enumerable by completeness of the fixed
effective calculus.  Finite validities therefore admit no effective calculus
that is both sound and complete for truth in all finite structures.

\section{Arithmetic formulas and hierarchy classes}
\label{sec:p04-arithmetic-hierarchy-syntax}

\begin{hbdefinition}{def:p04-language-arithmetic}{Language and standard structure of arithmetic}
The \emph{language of arithmetic} is
\[
 \Lang_A=\{0,S,+,\times\},
\]
where \(0\) is a constant, \(S\) is unary, and \(+\) and \(\times\) are
binary function symbols.  Its \emph{standard structure} is
\[
 \mathcal N_{\mathrm{ar}}=(\N,0,S,+,\times)
\]
with the usual operations.
\index{language of arithmetic}
\index{standard model of arithmetic}
\end{hbdefinition}

\begin{hbconvention}{def:p04-bounded-quantifiers}{Numerals, comparisons, and bounded quantifiers}
The numeral for \(n\in\N\) is \(\bar n=S^n0\); for
\(\vec a=(a_1,\ldots,a_k)\), put
\(\overline{\vec a}=(\bar a_1,\ldots,\bar a_k)\).  This bar is a
metalinguistic map from numbers to terms.

For arithmetic terms \(r,t\), write
\[
 r<t\;\defeq\;\exists z\,(r+Sz=t),
 \qquad
 r\leq t\;\defeq\;\exists z\,(z+r=t),
\]
using the least variable \(z\) absent from \(r,t\).  When
\(x\notin\operatorname{Var}(t)\), the marked forms
\[
\begin{aligned}
 \exists x<t\,\varphi&\defeq\exists x\,(x<t\land\varphi),&
 \forall x<t\,\varphi&\defeq\forall x\,(x<t\to\varphi),\\
 \exists x\leq t\,\varphi&\defeq\exists x\,(x\leq t\land\varphi),&
 \forall x\leq t\,\varphi&\defeq\forall x\,(x\leq t\to\varphi)
\end{aligned}
\]
count as bounded quantifiers before expansion.  In \(\mathcal N_{\mathrm{ar}}\),
\(x\leq t\) and \(x<St\) are equivalent.  No corresponding theorem of
Robinson arithmetic is asserted here.  The witness occurs to the left of
\(+\) in \(z+r=t\); this is the R-003 convention used by the bounded
comparison lemmas in Part~5.
\index{numeral}
\index{bounded quantifier}
\index{less than}
\index{less than or equal to}
\end{hbconvention}

\begin{hbdefinition}{def:p04-syntactic-hierarchy}{Syntactic arithmetical hierarchy}
A \emph{bounded arithmetic formula} is generated from atomic
\(\Lang_A\)-formulas using \(\bottom,\neg,\land,\lor,\to\) and marked
bounded quantifiers only.  Convert an arbitrary formula to negation normal
form by a fixed capture-avoiding procedure that eliminates implication,
pushes negation to atoms, applies bounded and unbounded quantifier dualities,
and never exchanges an unbounded quantifier with a bounded quantifier or a
connective.

The classes \(\Sigma^0_0\) and \(\Pi^0_0\) are exactly the bounded formulas
in negation normal form.  Having defined both classes at level \(n\), define
simultaneously:
\begin{itemize}
 \item \(\Sigma^0_{n+1}\) is the least class containing all
       \(\Sigma^0_n\)- and \(\Pi^0_n\)-formulas and closed under
       \(\land,\lor\), both marked bounded quantifiers, and unbounded
       existential quantification;
 \item \(\Pi^0_{n+1}\) is defined dually, using unbounded universal
       quantification.
\end{itemize}
An arbitrary formula belongs to one of these classes exactly when its fixed
negation-normal-form conversion does.  This is a syntactic classification,
not a classification by accidental truth in \(\mathcal N_{\mathrm{ar}}\).
\index{arithmetical hierarchy!syntactic}
\index{bounded arithmetic formula}
\end{hbdefinition}

\begin{hbdefinition}{def:p04-definability-classes}{Definability classes of relations}
For \(A\subseteq\N^k\), say that \(A\) is
\emph{\(\Sigma^0_n\)-definable} when some \(\Sigma^0_n\)-formula
\(\varphi(\vec x)\) satisfies
\[
 \vec a\in A\quad\Leftrightarrow\quad
 \mathcal N_{\mathrm{ar}}\entails\varphi(\overline{\vec a})
 \qquad(\vec a\in\N^k).
\]
Define \(\Pi^0_n\)-definability analogously and put
\[
\begin{aligned}
 \boldsymbol\Sigma^0_n(k)&=\{A\subseteq\N^k:A\text{ is }
   \Sigma^0_n\text{-definable}\},\\
 \boldsymbol\Pi^0_n(k)&=\{A\subseteq\N^k:A\text{ is }
   \Pi^0_n\text{-definable}\},\\
 \boldsymbol\Delta^0_n(k)&=
   \boldsymbol\Sigma^0_n(k)\cap\boldsymbol\Pi^0_n(k).
\end{aligned}
\]
The argument \(k=1\) may be omitted.  Unbold symbols classify formulas;
bold symbols classify relations definable in \(\mathcal N_{\mathrm{ar}}\).
\index{arithmetical hierarchy!definability classes}
\index{Delta-0-n definability@\(\boldsymbol\Delta^0_n\)-definability}
\end{hbdefinition}

\begin{hbconvention}{def:p04-convergence-notation}{Convergence notation}
For an expression \(F(\vec x)\) denoting a possibly partial numerical value,
write \(F(\vec x)\mathord\downarrow\) when that value is defined and
\(F(\vec x)\mathord\uparrow\) when it is undefined.  These are
metalinguistic predicates on partial computations, not numerical outputs.
The notation is needed to keep the domain and jump set-builders mathematical;
unlike Kleene equality, it classifies one expression rather than comparing
two expressions.
\index{convergence notation}
\index{divergence notation}
\end{hbconvention}

\section{Bounded coding and prefix form}
\label{sec:p04-bounded-coding}

\begin{hblemma}{thm:p04-pairing-graph}{proved=annex:p04-hierarchy.2}{Pairing graph}
Let
\[
 \mathsf P(x,y,z)\defeq(x+y)\times S(x+y)+(y+y)=z+z.
\]
Then \(\mathsf P\) is quantifier-free and, for \(a,b,c\in\N\),
\[
 \mathcal N_{\mathrm{ar}}\entails\mathsf P(\bar a,\bar b,\bar c)
 \quad\Leftrightarrow\quad
 c=\operatorname{pair}_{\N}(a,b).
\]
Moreover \(a,b\leq_{\N}\operatorname{pair}_{\N}(a,b)\).
\end{hblemma}

In the next result \(d\) abbreviates \(S((Si)\times m)\), and \(u\) is
fresh.  The final conjunct is deliberately the marked-bounded form of
\(v<d\), not the plain abbreviation whose expansion has an unbounded
existential.

\begin{hblemma}{thm:p04-goedel-beta}{proved=annex:p04-hierarchy.3}{G\"odel beta function}
Define
\[
 \mathsf B(w,i,v)\defeq
 \exists t<Sw\,\exists m<Sw\,\exists q<St\,
 \bigl(\mathsf P(t,m,w)\land t=d\times q+v\land
       \exists u<d\,(v+Su=d)\bigr).
\]
This is a bounded arithmetic formula.  For every \(w,i\in\N\) there is
exactly one \(v\leq_{\N} w\) satisfying \(\mathsf B(\bar w,\bar i,\bar v)\) in
\(\mathcal N_{\mathrm{ar}}\); write it \(\beta(w,i)\).  Conversely, every finite sequence
\(a_0,\ldots,a_{k-1}\) is \(\beta(w,\cdot)\) below \(k\) for some \(w\).
\end{hblemma}

The beta-function argument originates with G\"odel \cite{Godel1931}; see
also \cite[ch.~16]{BoolosBurgessJeffrey2007}.

\begin{hblemma}{thm:p04-pairing-contraction}{proved=annex:p04-hierarchy.4}{Pairing contraction}
Let \(z\) be distinct from \(x,y\) and not free in \(\varphi\).  Under every
assignment in \(\mathcal N_{\mathrm{ar}}\),
\[
 \exists x\,\exists y\,\varphi
 \quad\Leftrightarrow\quad
 \exists z\,\exists x<Sz\,\exists y<Sz\,
   (\mathsf P(x,y,z)\land\varphi),
\]
and dually
\[
 \forall x\,\forall y\,\varphi
 \quad\Leftrightarrow\quad
 \forall z\,\forall x<Sz\,\forall y<Sz\,
   (\mathsf P(x,y,z)\to\varphi).
\]
\end{hblemma}

\begin{hblemma}{thm:p04-bounded-exchange}{proved=annex:p04-hierarchy.5}{Bounded-quantifier exchange}
Let \(y,w\notin\operatorname{Var}(t)\), with \(w\) fresh.  Under every
assignment in \(\mathcal N_{\mathrm{ar}}\),
\[
 \forall y<t\,\exists x\,\varphi
 \quad\Leftrightarrow\quad
 \exists w\,\forall y<t\,\exists x<w\,\varphi,
\]
and
\[
 \exists y<t\,\forall x\,\varphi
 \quad\Leftrightarrow\quad
 \forall w\,\exists y<t\,\forall x<w\,\varphi.
\]
These are standard-structure equivalences used for boldface definability
classes; they are not derivability claims in a theory of arithmetic.
\end{hblemma}

\begin{hbdefinition}{def:p04-prefix-form}{Prefix-form hierarchy}
Two arithmetic formulas are \emph{\(\mathcal N_{\mathrm{ar}}\)-equivalent} when exactly
the same assignments satisfy them in \(\mathcal N_{\mathrm{ar}}\).  A formula is
\emph{prefix-form \(\Sigma^0_n\)} when it consists of \(n\) alternating
blocks of unbounded quantifiers, beginning existentially, applied to a bounded
formula in negation normal form.  Prefix-form \(\Pi^0_n\) begins universally.
\index{prefix normal form!arithmetical hierarchy}
\index{N-equivalence@\(\mathcal N_{\mathrm{ar}}\)-equivalence}
\end{hbdefinition}

\begin{hbtheorem}{thm:p04-prefix-normal-form}{proved=annex:p04-hierarchy.6}{Prefix normal form}
For every \(n\geq0\):
\begin{enumerate}
 \item relations defined by prefix-form \(\Sigma^0_n\)-formulas are closed,
       up to \(\mathcal N_{\mathrm{ar}}\)-equivalence of defining formulas, under
       \(\land,\lor\), both marked bounded quantifiers, and, for \(n\geq1\),
       unbounded existential quantification; dually for \(\Pi^0_n\);
 \item every \(\Sigma^0_n\)-formula is \(\mathcal N_{\mathrm{ar}}\)-equivalent to a
       prefix-form \(\Sigma^0_n\)-formula, and dually for \(\Pi^0_n\); and
 \item the boldface classes are exactly the relation classes defined in
       \(\mathcal N_{\mathrm{ar}}\) by prefix-form formulas of the corresponding kind.
\end{enumerate}
\end{hbtheorem}

\section{Computation histories}
\label{sec:p04-computation-histories}

\begin{hbdefinition}{def:p04-history-code}{History coding}
Use the fixed tuple coding to encode a finite configuration as its state and
head position, coding an integer \(h\) by \(2h\) when \(h\geq0\) and by
\(-2h-1\) when \(h<0\).  Code an everywhere-blank tape by \(0\).  Otherwise,
if \(\ell\) and \(r\) are its least and greatest nonblank positions, code the
tape by one plus the triple consisting of the signed codes of \(\ell,r\) and
the base-\(\abs{\mathcal A_{\mathrm{tape}}}\) numeral for
\(T(\ell),\ldots,T(r)\), most significant symbol first.  The endpoints supply
the block length, and every position outside the block decodes as blank.  This
separate zero case is necessary because an intermediate configuration may
have an entirely blank tape.  Use the program-numbering convention's symbol
numbers in the finite block, and encode finite lists of configuration tuples
by the beta function.  These choices determine one primitive-recursive
decoding, not a new machine model
\cite{Shoenfield1967,BoolosBurgessJeffrey2007}.

Read member \(j\) of entry family \(t\) from a single number \(s\) as
\(\beta(s,\operatorname{pair}_{\N}(t,j))\), using successive natural-number
tags for the families in (h1)--(h4), and successive subtags in their displayed
order for the component sequences within a clause.  A number \(s\) is a
\emph{history code}
for \((e,x)\) exactly when its decoded entries satisfy:
\begin{enumerate}
 \item[(h1)] the binary digits of \(x\), with certified partial values satisfying
       \(v_0=d_0\),
       \(v_{j+1}=\bar2\times v_j+d_{j+1}\) for every
       \(j+1<_{\N}L\), \(v_{L-1}=x\), and
       \(v_j\leq_{\N}x\) and \(d_j\leq_{\N}1\) for every
       \(j<_{\N}L\);
 \item[(h2)] the decoded program table is exactly the table determined by
       \(e\) under Convention~\ref{def:p04-program-numbering}, including all
       designated-slot validity checks;
 \item[(h3)] the decoded list is \(C_0,\ldots,C_T\), where \(C_0\) is the
       initial configuration on \(\operatorname{bin}(x)\) and
       \(C_j\longrightarrow_{\mathcal T_e}C_{j+1}\) for every
       \(j<_{\N}T\), with
       \(\mathcal T_e\) the decoded machine;
 \item[(h4)] a final accepting configuration whose prefix before the first blank is
       a valid binary numeral.
\end{enumerate}
Every auxiliary value is an entry of \(s\) or a component of a pairing code,
so the beta and pairing bounds control every quantifier introduced below.
\index{computation history code}
\end{hbdefinition}

\begin{hbtheorem}{thm:p04-bounded-verification}{sketched=annex:p04-hierarchy.7}{Bounded verification of computations}
There is a bounded \(\Lang_A\)-formula \(\psi(e,x,s)\) such that
\begin{enumerate}
 \item \(\mathcal N_{\mathrm{ar}}\entails\psi(\bar e,\bar x,\bar s)\) exactly
       when \(s\) is a history code for a valid accepting computation of
       program \(e\) on \(x\) with a valid numerical output; and
 \item \(\varphi_e(x)\) is defined exactly when
       \(\mathcal N_{\mathrm{ar}}\entails\exists s\,\psi(\bar e,\bar x,s)\).
\end{enumerate}
Order bounded \(\Lang_A\)-formula strings first by length and then
lexicographically by the symbol-registry order.  Define
\(\mathsf{Comp}(e,x,s)\) to be the first formula in that order satisfying the
two clauses.  The existence assertion makes this definition total;
\(\mathsf{Comp}\) names that arithmetic formula, not the metalinguistic
relation it defines, and a different adequate coding would generally choose a
different formula.
\end{hbtheorem}

\begin{hbtheorem}{thm:p04-ce-sigma-one}{sketched=annex:p04-hierarchy.8}{Computable enumerability and level one}
A set is computably enumerable exactly when it belongs to
\(\boldsymbol\Sigma^0_1\), and is decidable exactly when it belongs to
\(\boldsymbol\Delta^0_1\).
\end{hbtheorem}

\begin{hbdefinition}{def:p04-many-one-completeness}{Many-one hardness and completeness}
For a class \(\mathcal C\) of subsets of \(\N\), a set \(B\) is
\emph{\(\mathcal C\)-hard under many-one reductions} when
\(A\leq_m B\) for every \(A\in\mathcal C\), and is
\emph{\(\mathcal C\)-complete} when additionally \(B\in\mathcal C\).
This is a uniform lower-bound property relative to \(\mathcal C\) and
\(\leq_m\), not semantic completeness, deductive completeness, or proof
status.
\index{many-one completeness}
\index{many-one hardness}
\end{hbdefinition}

Let
\[
 K=\left\{e\in\N:\varphi_e(e)\mathord\downarrow\right\}
\]
be the self-application domain set of the preceding chapter.  Define the
bounded arithmetic formula
\[
 \mathsf{Acc}(e,x,k)\defeq
 \exists s\leq k\,\mathsf{Comp}(e,x,s).
\]
It asserts the existence of an accepting-history certificate whose code is at
most \(k\); it does not identify the certificate bound with a running-time
bound.  Put
\[
 \mathsf{TOT}=\left\{e\in\N:\forall x\,\exists k\,
                   \mathsf{Acc}(e,x,k)\right\}.
\]
The named relation keeps prose out of the set-builder and makes the displayed
\(\Pi^0_2\) upper bound literal.

\begin{hbtheorem}{thm:p04-domain-totality-completeness}{cited=Shoenfield1967}{Program domains and totality in the hierarchy}
The set \(K\) is \(\boldsymbol\Sigma^0_1\)-complete under many-one
reductions, its complement is \(\boldsymbol\Pi^0_1\)-complete, and
\(\mathsf{TOT}\) is \(\boldsymbol\Pi^0_2\)-complete.
\end{hbtheorem}

\begin{hbtheorem}{thm:p04-arithmetical-strictness}{cited=Shoenfield1967}{Strictness of the arithmetical hierarchy}
For every \(n\geq1\),
\[
 \boldsymbol\Delta^0_n\subsetneq\boldsymbol\Sigma^0_n,
 \qquad
 \boldsymbol\Delta^0_n\subsetneq\boldsymbol\Pi^0_n,
 \qquad
 \boldsymbol\Sigma^0_n\neq\boldsymbol\Pi^0_n,
\]
and, for unary relations,
\[
 \boldsymbol\Sigma^0_n\cup\boldsymbol\Pi^0_n
 \subsetneq\boldsymbol\Delta^0_{n+1}.
\]
\end{hbtheorem}

\begin{hbwarning}{def:p04-syntax-versus-class}{Syntactic prefix versus semantic class}
The syntactic level of a displayed formula gives an upper bound on the class
of the relation it defines.  A completeness reduction or nondefinability
argument is needed to exclude every lower class.
\end{hbwarning}

\begin{figure}[htbp]
\centering
\begin{tikzpicture}[
  node distance=11mm and 18mm,
  box/.style={draw,rounded corners,minimum width=22mm,minimum height=7mm,
              align=center,font=\small},
  inc/.style={-{Latex[length=2mm]},thick}
]
 \node[box] (dn) {\(\boldsymbol\Delta^0_n\)};
 \node[box,above left=of dn] (sn) {\(\boldsymbol\Sigma^0_n\)};
 \node[box,above right=of dn] (pn) {\(\boldsymbol\Pi^0_n\)};
 \node[box,above=24mm of dn] (dnext) {\(\boldsymbol\Delta^0_{n+1}\)};
 \node[box,above left=of dnext] (snext) {\(\boldsymbol\Sigma^0_{n+1}\)};
 \node[box,above right=of dnext] (pnext) {\(\boldsymbol\Pi^0_{n+1}\)};
 \draw[inc] (dn) -- node[left,font=\scriptsize] {strict} (sn);
 \draw[inc] (dn) -- node[right,font=\scriptsize] {strict} (pn);
 \draw[inc] (sn) -- node[left,font=\scriptsize] {strict} (dnext);
 \draw[inc] (pn) -- node[right,font=\scriptsize] {strict} (dnext);
 \draw[inc] (dnext) -- node[left,font=\scriptsize] {strict} (snext);
 \draw[inc] (dnext) -- node[right,font=\scriptsize] {strict} (pnext);
 \node[font=\scriptsize,below=2mm of dn] {\(n\geq1\)};
\end{tikzpicture}
\caption{Arithmetical-hierarchy inclusions (apparatus A-124).  Separate
upper nodes record \(\boldsymbol\Sigma^0_n\neq\boldsymbol\Pi^0_n\);
arrows marked strict are exactly those asserted by
Theorem~\ref{thm:p04-arithmetical-strictness}.}
\label{fig:p04-arithmetical-hierarchy}
\end{figure}

%% Migration accounting for v1.1 lines 8248--8897:
%% - every definition, warning, and theorem-class result is retained;
%% - all source proof statuses are preserved, including the formerly unlabeled
%%   but fully proved complement-domain corollary's upstream use in the table;
%% - exercises at 8877--8896 are deliberately retired to verification assets:
%%   they add practice prompts but no new definitions or result statements.
